\chapter{Zusammenfassung}
\label{chap:zusammenfassung}

In der Arbeit wurde untersucht, warum sich Kryptowährungen bisher nicht als alltägliches Zahlungsmittel durchgesetzt haben. Ein Vergleich von El Salvador, den Vereinigten Arabischen Emiraten und der Europäischen Union zeigt, dass weder eine gesetzliche Einführung noch eine staatliche Regulierung allein ausreicht. In allen drei Fällen blieb die tatsächliche Nutzung im Alltag gering, obwohl die Ausgangsbedingungen und Strategien unterschiedlich waren.\\

Die zentrale Ursache dafür ist, dass Kryptowährungen die klassischen Geldfunktionen – Tauschmittel, Recheneinheit und Wertaufbewahrungsmittel – aufgrund von Preisvolatilität, mangelndem Vertrauen und fehlendem Alltagsnutzen nur eingeschränkt erfüllen. Für eine breitere Akzeptanz wäre das Zusammenspiel dieser drei Faktoren entscheidend.\\

Offene Fragen betreffen vor allem die langfristige Entwicklung staatlich kontrollierter digitaler Währungen wie des digitalen Euros und deren Wirkung auf das gesellschaftliche Vertrauen in digitale Zahlungsmittel.